\documentclass{article}
\usepackage{hyperref}
\usepackage{xcolor}

\hypersetup{colorlinks,linkcolor=blue,urlcolor=blue}

\title{Why We Rewrote Our Auth Layer (A Post-Mortem)}
\author{Engineering Blog}
\date{April 2026}

\begin{document}
\maketitle

If you're here, you've probably seen the incident report already. This post covers what led to the rewrite, what we tried, and what worked.

\section*{Background}
Our auth layer was written five years ago when we had fewer than 10{,}000 users. It aged out around the 1M mark --- token validation alone was adding 40ms to every request.

\section*{What we tried first}
We tried to optimise the existing implementation. Moved validation to a cache. Cut the per-request overhead in half. Not enough.

Then we tried sharding the session store. Helped, but the original design was never meant for this scale, and we kept hitting edge cases.

\section*{The rewrite}
After about a week of trying to shore up the old system, we made the call to rewrite. Three engineers, six weeks, a parallel rollout with shadow traffic comparison.

Key design choices:
\begin{itemize}
  \item Stateless tokens (JWT) instead of session-backed lookups
  \item Revocation list cached in Redis with 30-second TTL
  \item Asymmetric keys so verify doesn't require a shared secret
\end{itemize}

\section*{Rollout}
Shadow traffic for two weeks, then a 1\% canary, then 10\%, then 100\%. No incidents.

\section*{Results}
Token validation now adds 2ms per request (was 40ms). The old auth service has been retired.

\section*{Lessons}
Try to optimise before rewriting. When you do rewrite, run in shadow long enough to catch edge cases. Rollout gradually even when shadow looks clean.

Questions? Email us at \href{mailto:eng@example.com}{eng@example.com}.

\end{document}
